% !TEX program = pdflatex
% Compile on Overleaf (free, no install) — https://overleaf.com
% Or locally: `pdflatex Greenify-Report.tex` (twice, for cross-refs)
% Or with Tectonic:                `tectonic Greenify-Report.tex`
\documentclass[11pt,a4paper]{article}

% ---------- Packages ----------
\usepackage[utf8]{inputenc}
\usepackage[T1]{fontenc}
\usepackage[margin=1in]{geometry}
\usepackage{graphicx}
\usepackage{xcolor}
\usepackage{hyperref}
\usepackage{fancyhdr}
\usepackage{titlesec}
\usepackage{booktabs}
\usepackage{tabularx}
\usepackage{enumitem}
\usepackage{listings}
\usepackage{tcolorbox}
\usepackage{microtype}
\usepackage{parskip}
\usepackage{float}

% ---------- Colours ----------
\definecolor{ink}{HTML}{0A0A0A}
\definecolor{accent}{HTML}{C0392B}
\definecolor{muted}{HTML}{555555}
\definecolor{code-bg}{HTML}{F5F2EC}

% ---------- Hyperlinks ----------
\hypersetup{
    colorlinks=true,
    linkcolor=accent,
    urlcolor=accent,
    citecolor=accent,
    pdftitle={Greenify — Technical Report},
    pdfauthor={Rajdeep Chaudhari},
    pdfsubject={CN6035 Mobile and Distributed Systems — Task 1}
}

% ---------- Sections ----------
\titleformat{\section}{\Large\bfseries\color{ink}}{\thesection.}{0.5em}{}
\titleformat{\subsection}{\large\bfseries\color{ink}}{\thesubsection}{0.5em}{}

% ---------- Header / footer ----------
\pagestyle{fancy}
\fancyhf{}
\fancyhead[L]{\small\color{muted}Greenify — CN6035 Task 1}
\fancyhead[R]{\small\color{muted}Rajdeep Chaudhari}
\fancyfoot[C]{\small\thepage}
\renewcommand{\headrulewidth}{0.4pt}

% ---------- Code listing style ----------
\lstset{
    backgroundcolor=\color{code-bg},
    basicstyle=\ttfamily\footnotesize,
    keywordstyle=\color{accent}\bfseries,
    commentstyle=\color{muted}\itshape,
    stringstyle=\color{accent},
    numbers=none,
    frame=single,
    rulecolor=\color{muted!40},
    breaklines=true,
    showstringspaces=false,
    captionpos=b,
    xleftmargin=0.4em,
    xrightmargin=0.4em
}

% ---------- Title block ----------
\title{\vspace{-1em}\textbf{Greenify}\\[0.3em]
{\large A hybrid DApp for on-chain carbon credits}\\[1em]
{\normalsize CN6035 Mobile and Distributed Systems — Task 1 (70\%)}}
\author{Rajdeep Chaudhari \quad\textbullet\quad Student ID 2522821\\
\href{https://github.com/rajdeepchaudhari-work/Greenify}{github.com/rajdeepchaudhari-work/Greenify}
\quad\textbullet\quad
\href{https://greenifyrc.vercel.app}{greenifyrc.vercel.app}}
\date{April 2026}

\begin{document}
\maketitle
\thispagestyle{fancy}

\begin{tcolorbox}[colback=code-bg, colframe=ink, boxrule=0.8pt, arc=0pt,
    left=8pt, right=8pt, top=6pt, bottom=6pt]
\textbf{Abstract.}
Greenify re-implements the carbon credit ledger on Ethereum. Three audited
smart contracts issue, transfer and retire ERC-1155 tokens where one unit
equals one tonne of CO\textsubscript{2} equivalent. A Vite + React SPA
brokers MetaMask interactions; a Vercel serverless API caches on-chain
events into MongoDB for fast reads and pins project metadata to IPFS via
Pinata. Contracts are deployed and source-verified on Sepolia. The project
hits 91\% Hardhat statement coverage, enforces lint+format+test gates in
CI, and ships a public-browsing mode that requires no wallet for read-only
interaction.
\end{tcolorbox}

\vspace{0.6em}

% ===========================================================
\section{Why I built this}
The voluntary carbon market moved roughly \$900\,bn last year. Most of it
still runs on PDFs, spreadsheets and proprietary registries you have to
email to audit. In 2023 a \emph{Guardian} investigation of Verra --- the
biggest certifier --- found the majority of its rainforest credits did not
represent real emission reductions. That isn't a Verra problem. It's a
ledger problem. When the record is private, double-counting, silent
re-use and retirement fraud are invisible to the buyers paying for the
offsets.

Greenify moves that record onto Ethereum. Three small contracts. Every
registration, mint, sale and retirement is a public transaction. Anyone
can audit the lifecycle of any credit from a block explorer. The hybrid
DApp on top is a nicer way to interact, but the ledger is the thing that
matters. This report walks through what I built, why I made the calls I
did, and where I'd take it next.

% ===========================================================
\section{System overview}

\begin{tcolorbox}[colback=code-bg, colframe=ink!80, boxrule=0.6pt, arc=0pt]
\begin{verbatim}
React + Vite SPA  ─►  Vercel serverless /api/*  ─►  MongoDB Atlas
     │                       │                        │
     │ ethers v6             │ ethers v6              │ cached events
     ▼                       ▼                        ▲
 MetaMask ──── signed tx ──►  Ethereum Sepolia ── logs ┘
                             ┌───────────────────────┐
                             │ ProjectRegistry       │
                             │ CarbonCredit (ERC-1155)│
                             │ Marketplace           │
                             └───────────────────────┘
\end{verbatim}
\end{tcolorbox}

Three tiers, each with a single job. Contracts own state. The serverless
API caches events and pins metadata. The SPA presents the ledger and
talks to the wallet. I kept the boundary between on-chain and off-chain
deliberately strict: Mongo is a \emph{read-only mirror} and can be wiped
and rebuilt from Sepolia at any time. No off-chain service is ever the
source of truth.

\begin{table}[H]
\centering
\small
\begin{tabularx}{\linewidth}{lX}
\toprule
\textbf{Layer} & \textbf{Stack} \\
\midrule
Smart contracts & Solidity 0.8.24, OpenZeppelin 5.0, Hardhat, TypeChain, Chai \\
Serverless API  & \texttt{@vercel/node}, Mongoose, Zod, Formidable, pinata-web3 \\
Frontend        & React 18, Vite 5, TailwindCSS 3, ethers v6, React Router 6 \\
Infrastructure  & Vercel (SPA + serverless), MongoDB Atlas M0, Pinata IPFS \\
Quality         & ESLint, Solhint, Prettier, Husky, lint-staged, GitHub Actions \\
\bottomrule
\end{tabularx}
\caption{The complete stack.}
\end{table}

% ===========================================================
\section{Smart contract design}

\subsection{ProjectRegistry}
The registry is the list of environmental projects.
\texttt{registerProject(string ipfsCid)} is permissionless; anyone can
submit a project and pin its metadata to IPFS. What it cannot do on its
own is issue credits. That gate is \texttt{VERIFIER\_ROLE} from
OpenZeppelin's \texttt{AccessControl}. For the coursework deployment the
verifier is a single test account; in production it would be a Gnosis
Safe multisig. I chose role-based access over plain \texttt{Ownable} so
that swap becomes a configuration change, not a contract rewrite.

One small detail I care about: the project id is the same number as the
ERC-1155 \texttt{tokenId} of the credit. There is no mapping to keep in
sync between the registry and the token. You register a project, you get
an id, credits minted against that id are that project --- done.

\subsection{CarbonCredit (ERC-1155)}
Credits are ERC-1155 tokens where one unit of any \texttt{tokenId}
equals one tonne of CO\textsubscript{2} equivalent. I looked at the
alternatives before picking ERC-1155:

\begin{itemize}[leftmargin=1.3em, itemsep=2pt]
    \item \textbf{ERC-20} --- collapses every credit into one fungible
    pool. You lose the ability to say ``I retired credits from \emph{this}
    forestry project in Brazil.'' Deal-breaker for carbon accounting.
    \item \textbf{ERC-721} --- every credit is a unique NFT. Way more
    expensive in gas, and over-modelled: credits inside a project
    \emph{are} fungible with each other.
    \item \textbf{ERC-1155} --- one contract, many project batches, each
    fungible internally. The right shape.
\end{itemize}

\texttt{mint()} is behind \texttt{MINTER\_ROLE} and calls
\texttt{registry.recordIssuance()} so the registry's issuance tally stays
consistent. The function that matters most is \texttt{retire()}. It
burns the caller's tokens and emits \texttt{CreditsRetired}. That event
is the receipt. Auditors, regulators, ESG reporters --- all of them can
resolve a retirement claim to a specific burn transaction on Sepolia. No
recycling, no re-listing, no second life.

\subsection{Marketplace}
An escrowed ERC-1155 listing book. Sellers call \texttt{list()}, which
immediately transfers their credits into the marketplace contract
(one-time \texttt{setApprovalForAll}). Buyers call \texttt{buy()} with
ETH. The contract transfers credits, pays the seller, takes a protocol
fee, and refunds any overpayment --- all in one atomic call.

\begin{lstlisting}[language={},caption={Buy flow (simplified) --- checks-effects-interactions in practice.}]
function buy(uint256 id, uint256 amount) external payable nonReentrant {
    Listing storage l = listings[id];
    if (!l.active)                             revert ListingInactive();
    if (amount == 0 || amount > l.amount)       revert NotEnoughAvailable();

    uint256 total = amount * l.pricePerUnit;
    if (msg.value < total)                      revert InsufficientPayment();

    uint256 fee            = (total * feeBps) / 10_000;
    uint256 sellerProceeds = total - fee;

    // effects BEFORE external calls
    l.amount -= amount;
    if (l.amount == 0) l.active = false;

    // interactions
    credits.safeTransferFrom(address(this), msg.sender, l.projectId, amount, "");
    _payout(l.seller, sellerProceeds);
    if (fee > 0)            _payout(feeRecipient, fee);
    if (msg.value > total)  _payout(msg.sender, msg.value - total);

    emit Sold(id, msg.sender, amount, total);
}
\end{lstlisting}

Four things I'm glad I nailed here:

\begin{enumerate}[leftmargin=1.3em, itemsep=2pt]
    \item \textbf{Checks~$\rightarrow$~effects~$\rightarrow$~interactions.}
    State is updated before any external call. Textbook re-entrancy
    defence.
    \item \textbf{\texttt{ReentrancyGuard} anyway.} Belt and braces. ETH
    transfers use low-level \texttt{call} with no gas limit, so in
    principle a malicious seller contract could try to re-enter. I'm not
    going to rely on ordering alone.
    \item \textbf{Custom errors everywhere.} \texttt{NotSeller},
    \texttt{InsufficientPayment}, \texttt{InvalidFee}. Cheaper than
    string reverts and far easier to decode on the client.
    \item \textbf{Fee capped in two places.} 1000 bps (10\%) hard-cap,
    enforced in the constructor \emph{and} in \texttt{setFee}. An owner
    cannot accidentally or deliberately set a 50\% fee.
\end{enumerate}

All three contracts are \textbf{source-verified on Etherscan}. The
bytecode on Sepolia matches the source in this repository, byte for
byte.

% ===========================================================
\section{Back-end: the indexer problem}
This is the part I most enjoyed re-architecting. The first version was a
long-running Express server with ethers' WebSocket event listeners. It
worked locally; it fell apart on free-tier hosting. Render's free dyno
sleeps after fifteen minutes of idle --- exactly long enough for public
RPCs to drop the subscription. The server wakes up, the subscription is
dead, and events start going to \texttt{/dev/null} silently. Not great.

I rewrote the indexer as a \textbf{checkpointed poll} and moved
everything into Vercel serverless functions colocated with the frontend.
One deploy, one URL, no Express.

\texttt{/api/sync} works like this:

\begin{itemize}[leftmargin=1.3em, itemsep=2pt]
    \item A single \texttt{Meta} document in Mongo stores
    \texttt{lastBlock} --- the last block the indexer processed.
    \item The endpoint reads that, calls
    \texttt{registry.queryFilter(...)} and
    \texttt{market.queryFilter(...)} in 45,000-block chunks (public RPCs
    cap \texttt{eth\_getLogs} at 50k per call), writes each event as a
    Mongo upsert, and advances the checkpoint.
    \item Vercel cron fires it daily (the free-tier minimum); a
    secondary cron-job.org monitor hits the endpoint every minute for
    sub-hour freshness.
\end{itemize}

The nice property is that the design is \textbf{idempotent by
construction}. Every write is an \texttt{updateOne} keyed by the event's
natural id (\texttt{projectId}, \texttt{listingId}). Run it twice on the
same range --- no change. Wipe Mongo entirely --- the next sync
rebuilds the whole state from Sepolia.

Other serverless-specific things I had to deal with:

\begin{itemize}[leftmargin=1.3em, itemsep=2pt]
    \item \textbf{Connection pooling.} Mongoose is expensive to
    initialise. I cache the connection promise on
    \texttt{global.\_mongoPromise} so warm invocations reuse the pool.
    \item \textbf{Multipart uploads.} Vercel's Node runtime skips
    body-parsing for \texttt{multipart/form-data}, which is what I needed
    for image uploads. \texttt{formidable} reads the raw stream; images
    are capped at 4\,MB.
    \item \textbf{ESM extensions.} Vercel runs compiled
    \texttt{api/*.js} under Node's ESM loader, which requires explicit
    \texttt{.js} extensions on relative imports. TypeScript lets you
    write \texttt{from './mongo'} but the runtime wants
    \texttt{from './mongo.js'}. Got bitten by this in production before
    I understood why.
\end{itemize}

% ===========================================================
\section{Frontend}
React + Vite. Two route trees: \texttt{/} is the marketing landing
(public), \texttt{/app/*} is the application behind a navbar with
Dashboard / Registry / Market.

State is small. Two contexts:

\begin{itemize}[leftmargin=1.3em, itemsep=2pt]
    \item \texttt{WalletProvider} wraps MetaMask --- exposes address,
    chainId, signer, and a \texttt{connect()} action. Also detects when
    no browser wallet is present, in which case write actions show an
    \texttt{InstallWalletCard} linking to MetaMask, Rabby and Coinbase
    Wallet instead of silently failing.
    \item \texttt{TxProvider} is a four-state banner:
    \texttt{signing $\rightarrow$ mining $\rightarrow$ success | error}.
    Every on-chain write goes through
    \texttt{tx.run("Mint credits", () => contract.mint(...))}. The
    banner drives itself, shows the tx hash as an Etherscan link, and
    catches \emph{user rejected} / \emph{insufficient funds} into
    friendly copy.
\end{itemize}

Three deliberate product choices:

\begin{enumerate}[leftmargin=1.3em, itemsep=2pt]
    \item \textbf{Public by default.} The read API has no auth. Someone
    landing on the site can see every project, listing and stat without
    ever clicking Connect Wallet. The wallet is only demanded when the
    user actually wants to sign something.
    \item \textbf{IPFS metadata in the client.} Project name,
    description, and image come from Pinata, fetched client-side with
    \texttt{ipfs.io} and Cloudflare as fallbacks.
    \item \textbf{Neo-brutalist design.} 3\,px black borders, hard
    offset shadows, Bricolage Grotesque on JetBrains Mono, high-contrast
    cream/red/yellow/green palette. Loud on purpose, accidentally
    accessible (contrast ratios and focus outlines satisfy WCAG AA
    without extra work).
\end{enumerate}

A top-level \texttt{ErrorBoundary} catches anything uncaught and renders
a recover-and-retry screen instead of a blank page.

% ===========================================================
\section{Quality and CI}
The repo is an npm workspace with two member packages
(\texttt{contracts}, \texttt{frontend}). GitHub Actions runs on every
push: ESLint with zero warnings allowed, Solhint, Prettier format-check,
and the full Hardhat test suite. Husky + lint-staged replay the same
checks on every local commit.

\begin{table}[H]
\centering
\small
\begin{tabularx}{\linewidth}{lrrrr}
\toprule
\textbf{File} & \textbf{\% Stmts} & \textbf{\% Branch} & \textbf{\% Funcs} & \textbf{\% Lines} \\
\midrule
\texttt{ProjectRegistry.sol} & 93.75 & 71.43 & 85.71 & 95.00 \\
\texttt{CarbonCredit.sol}    & 81.82 & 62.50 & 60.00 & 83.33 \\
\texttt{Marketplace.sol}     & 92.86 & 56.67 & 100.00 & 91.67 \\
\midrule
\textbf{All files}           & \textbf{90.91} & \textbf{61.54} & \textbf{83.33} & \textbf{91.18} \\
\bottomrule
\end{tabularx}
\caption{Hardhat coverage (13 passing tests).}
\end{table}

Everything below is green:
\begin{itemize}[leftmargin=1.3em, itemsep=2pt]
    \item 13 passing Chai tests covering the full
    register$\rightarrow$approve$\rightarrow$mint$\rightarrow$list$\rightarrow$buy$\rightarrow$cancel$\rightarrow$retire flow.
    \item Solhint: zero violations, pinned to Solidity \textasciicircum 0.8.24.
    \item ESLint: zero warnings with \texttt{@typescript-eslint} +
    \texttt{react-hooks} strict config.
    \item TypeScript strict mode across the Vite project and the
    separate \texttt{tsconfig.node.json} for the serverless API.
\end{itemize}

Commits use conventional-commits prefixes and are attributed to my
GitHub account via a noreply email.

% ===========================================================
\section{Security --- what I thought about}
I didn't run a formal audit (this is coursework), but I walked through
the obvious threats and documented how each is mitigated.

\begin{table}[H]
\centering
\small
\begin{tabularx}{\linewidth}{lX}
\toprule
\textbf{Threat} & \textbf{Mitigation} \\
\midrule
Re-entrancy on \texttt{buy}   & Checks-effects-interactions + \texttt{ReentrancyGuard} \\
Unauthorised minting          & \texttt{MINTER\_ROLE}; mint also checks registry approval \\
Double-counting               & ERC-1155 with unique project ids; transfers move credits, retire burns them \\
Retirement fraud              & \texttt{retire()} is a real burn; \texttt{CreditsRetired} event is the audit trail \\
Owner raising fees mid-market & 10\% cap in constructor \emph{and} \texttt{setFee} \\
Silent ETH payout failure     & Every \texttt{.call} is return-checked with \texttt{TransferFailed} \\
Leaked secrets                & \texttt{.env} gitignored; Vercel holds secrets encrypted \\
Malicious verifier            & Role revocable by admin; production needs a multisig \\
\bottomrule
\end{tabularx}
\caption{Threat model.}
\end{table}

For mainnet I would add Slither, Mythril, and a Foundry invariant test
suite before going live. None are hard; they take time I didn't have
this sprint.

% ===========================================================
\section{What's not here yet}
Honest limitations I'd address in v2:

\begin{itemize}[leftmargin=1.3em, itemsep=2pt]
    \item \textbf{Verifier centralisation.} A single EOA holds
    \texttt{VERIFIER\_ROLE}. Fine for a test deployment, fatal in
    production. Swap for a Gnosis Safe with a formal MRV
    (measurement-reporting-verification) pipeline behind it.
    \item \textbf{Cron cadence.} Vercel's free tier caps cron at daily.
    cron-job.org fills the gap, but a paid plan or self-hosted trigger
    is the proper long-term answer.
    \item \textbf{Chain lock-in.} Sepolia only. The contracts are
    chain-agnostic --- adding Polygon or Arbitrum is a config change plus
    a small chain-switcher in the UI. No Solidity changes required.
    \item \textbf{No off-chain attestation trail.} Credits are backed by
    whatever the verifier looked at off-chain. An EIP-712-signed
    attestation document would close that loop cleanly. Natural next
    feature.
\end{itemize}

% ===========================================================
\section{Deliverables}

\begin{table}[H]
\centering
\small
\begin{tabularx}{\linewidth}{lX}
\toprule
\textbf{Deliverable} & \textbf{Location} \\
\midrule
Source code (contracts + frontend) & \texttt{/contracts}, \texttt{/frontend} \\
Tests (13 passing, 91\% cov)        & \texttt{contracts/test/Greenify.test.ts} \\
CI pipeline                         & \texttt{.github/workflows/ci.yml} \\
Deployed addresses                  & \texttt{contracts/deployments/sepolia.json} \\
Installation manual                 & \texttt{SETUP.md} \\
Repo README                         & \texttt{README.md} \\
Live demo                           & \url{https://greenifyrc.vercel.app} \\
\bottomrule
\end{tabularx}
\end{table}

\subsection*{Deployed contract addresses (Sepolia)}
\begin{table}[H]
\centering
\small
\begin{tabularx}{\linewidth}{lX}
\toprule
\textbf{Contract} & \textbf{Address} \\
\midrule
ProjectRegistry          & \href{https://sepolia.etherscan.io/address/0x86B861a6F7E4B10CD96B7491fFA6a0967441142b\#code}{\texttt{0x86B861a6F7E4B10CD96B7491fFA6a0967441142b}} \\
CarbonCredit (ERC-1155)  & \href{https://sepolia.etherscan.io/address/0x4272be407BA26d9aD469E5951e549eD36932bB9E\#code}{\texttt{0x4272be407BA26d9aD469E5951e549eD36932bB9E}} \\
Marketplace              & \href{https://sepolia.etherscan.io/address/0x8BEdAf9e29aC4FBa25D679Dc8Fa4AdAc126a6403\#code}{\texttt{0x8BEdAf9e29aC4FBa25D679Dc8Fa4AdAc126a6403}} \\
\bottomrule
\end{tabularx}
\end{table}

% ===========================================================
\section{References}

\begin{enumerate}[leftmargin=1.3em, itemsep=2pt]
    \item UNDP Climate Promise (2022). \emph{What are carbon markets and
    why are they important?}
    \url{https://climatepromise.undp.org/news-and-stories/what-are-carbon-markets-and-why-are-they-important}
    \item The Guardian (2023). \emph{Revealed: more than 90\% of
    rainforest carbon offsets by biggest certifier are worthless,
    analysis shows.}
    \url{https://www.theguardian.com/environment/2023/jan/18/revealed-forest-carbon-offsets-biggest-provider-worthless-verra-aoe}
    \item World Bank (2023). \emph{State and Trends of Carbon Pricing
    2023.}
    \url{https://openknowledge.worldbank.org/entities/publication/58f2a409-9bb7-4ee6-899d-be47835c838f}
    \item OpenZeppelin Contracts v5.0 documentation.
    \url{https://docs.openzeppelin.com/contracts/5.x/}
    \item Ethereum EIP-1155: Multi Token Standard.
    \url{https://eips.ethereum.org/EIPS/eip-1155}
    \item Vercel Serverless Functions --- Node.js runtime docs.
    \url{https://vercel.com/docs/functions/runtimes/node-js}
    \item MongoDB (2024). \emph{Atlas free cluster connection best
    practices.}
    \url{https://www.mongodb.com/docs/atlas/reference/free-shared-limitations/}
\end{enumerate}

\end{document}
